\documentclass[aspectratio=169,t]{beamer} % widescreen 16:9; top-align frame content
% If you want widescreen use: \documentclass[aspectratio=169]{beamer}

% --- Colors ---
\usepackage{xcolor}
\definecolor{mygreen}{RGB}{109, 141, 36}
\definecolor{myblue}{RGB}{34,31,217}
\definecolor{mybrown}{RGB}{194,164,113}
\definecolor{myred}{RGB}{255,66,56}

% --- Fonts/encoding ---
\usepackage[T1]{fontenc}
\usepackage{lmodern}
\usefonttheme{professionalfonts}

% --- Packages ---
\usepackage{amsmath,amsthm,bm}
\usepackage{graphicx}
\usepackage{booktabs,multicol}
\usepackage{placeins}
\usepackage{tikz}
\usepackage{todonotes}
\usetikzlibrary{positioning}
\usepackage[english]{babel}
\usepackage{microtype}
\usepackage{apacite}
\presetkeys{todonotes}{disable}{}

\makeatletter

\newcommand{\safeincludegraphics}[2][]{%
  \IfFileExists{#2}{%
    \includegraphics[#1]{#2}%
  }{%
    \fbox{%
      \parbox[c][0.20\textheight][c]{0.9\linewidth}{%
        \centering Missing graphic\\\texttt{\detokenize{#2}}%
      }%
    }%
  }%
}

\makeatother

\newcommand{\safeinput}[1]{%
  \IfFileExists{#1}{%
    \begin{center}
    \fbox{%
      \parbox{0.9\linewidth}{%
        \centering External table omitted in compile-safe slide mode\\\texttt{\detokenize{#1}}%
      }%
    }
    \end{center}
  }{%
    \begin{center}
    \fbox{%
      \parbox{0.9\linewidth}{%
        \centering Missing input file\\\texttt{\detokenize{#1}}%
      }%
    }
    \end{center}
  }%
}

% --- Presentation info ---
%\title[Forecasting S\&P 500 Direction]{Forecasting S\&P 500 Direction\\[0.3em]
%\large Historical Signals and Market Direction}

\title[Historical Signals and Market Direction]{Historical Signals and Market Direction}
\author{Trang Nguyen} 
\institute{STAT 587: Data Science \\ Portland State University}
\date{March 9th, 2026}

% --- Remove default navigation symbols ---
\setbeamertemplate{navigation symbols}{}
\setbeamertemplate{caption}[numbered]
\setcounter{tocdepth}{2}

% --- Custom header (headline) ---
\setbeamertemplate{headline}{%
  \begin{tikzpicture}[remember picture,overlay]
    % header bar height tuned for better top padding
    \fill[mygreen] (current page.north west) rectangle ([yshift=-0.46cm]current page.north east);

    \node[anchor=west, text=white, font=\bfseries\small, xshift=0.35cm, yshift=-0.28cm]
      at (current page.north west) {Historical Signals and Market Direction};
  \end{tikzpicture}
  \vspace{0.35cm} % reserve room under custom header to avoid overlap
}

% --- Global frame title style for content slides ---
\setbeamerfont{frametitle}{size=\large,series=\bfseries}
\setbeamercolor{frametitle}{fg=black,bg=white}
\setbeamertemplate{frametitle}{%
  \vspace{0.20cm}
  \begin{beamercolorbox}[wd=\paperwidth,leftskip=0.6cm,rightskip=0.6cm,sep=0.25ex]{frametitle}
    \usebeamerfont{frametitle}\insertframetitle
  \end{beamercolorbox}
  \vspace{0.10cm}
}

\setbeamertemplate{footline}{%
  \begin{tikzpicture}[remember picture,overlay]
    % thicker footer bar (adjust 0.55cm)
    \fill[mygreen] (current page.south west) rectangle ([yshift=0.40cm]current page.south east);

    % footer text placed ON the green bar
    \node[anchor=west, text=white, font=\scriptsize, xshift=0.35cm, yshift=0.18cm]
      at (current page.south west)
      {\insertshortauthor\ \raisebox{0.2mm}{$\bm{\vert}$}\ PSU};

    \node[anchor=east, text=white, font=\scriptsize, xshift=-0.35cm, yshift=0.18cm]
      at (current page.south east)
      {\insertframenumber/\inserttotalframenumber};
  \end{tikzpicture}
  \vspace{0.05cm}
}


% --- Optional: a simple highlight box command ---
\newcommand{\mybox}[1]{%
  \begin{beamercolorbox}[rounded=true,shadow=false,sep=6pt,wd=\linewidth]{block body}
    \color{black}#1
  \end{beamercolorbox}
}


\begin{document}

% --- Title slide ---
\begin{frame}[plain]
\begin{tikzpicture}[remember picture,overlay]
  % ---- Rectangle size (edit these) ----
  \def\BoxW{\paperwidth}
  \def\BoxH{0.45\paperheight}

  % ---- Green rectangle centered on the page ----
  \fill[mygreen] 
    ([xshift=-0.5*\BoxW,yshift=-0.5*\BoxH]current page.center)
    rectangle
    ([xshift= 0.5*\BoxW,yshift= 0.5*\BoxH]current page.center);

  % ---- Centered text within the rectangle ----
  \node[
    align=center,
    text=white,
    font=\sffamily,
    text width=0.75\paperwidth
  ] at (current page.center) {%
    {\Large\bfseries \inserttitle\par}
    \vspace{0.4cm}
    {\normalsize \insertauthor\par}
    {\normalsize \insertdate\par}
  };
\end{tikzpicture}
\end{frame}





% --- Outline slide ---
\begin{frame}{Outline}
  \tableofcontents
\end{frame}


\section{Introduction}

\subsection{Motivation}

\begin{frame}{Motivation}
\begin{itemize}
    \item Market direction forecast is important for trading decisions.
    \item Predictability reveals whether historical signals are informative about market direction.
    \item It is an empirical test of weak-form of the \textbf{Efficient Market Hypothesis (EMH)} in practice.
    \item If historical signals are predictive out of sample, this challenges weak-form EMH.
    \item If not, the results are more consistent with the \textbf{EMH}.

\end{itemize}
\end{frame}

\subsection{Research Question}
\begin{frame}{Research Question}
\begin{itemize}
  \item Objective: Forecasting whether the S\&P 500 goes up or down the next trading day.
  \item Main question: Do historical price and volume contain usable predictive information?
  \item Scope: Prediction from market history only, without macro or news variables.
\end{itemize}
\end{frame}
\begin{frame}{Modeling Strategy}
\begin{itemize}
  \item Engineer features from OHLCV data.
  \item Reduce dimensionality through screening and selection.
  \item Fit classification models aligned with the data structure and characteristics.
  \item Tune hyperparameters using iterative cross validation evaluation.
\end{itemize}
\end{frame}


\subsection{Background}
\begin{frame}{Background}
\begin{itemize}
  \item Weak-form EMH: past prices and volume should not predict future returns.
  \item S\&P 500: primary large-cap U.S. equity benchmark.
  \item Forecasting challenge: many firm-level signals plus shifting, sometimes abrupt, political and macroeconomic conditions.
  \item Study focus: whether past market signals retain predictive value in today’s faster and rapidly updating market.
\end{itemize}
\end{frame}






\section{Dataset \& preprocessing}
\subsection{Dataset \& preprocessing}
\begin{frame}{Dataset}
\begin{itemize}
  \item Source: Yahoo Finance via \texttt{yfinance}.
  \item Data cover S\&P 500 stocks and the S\&P 500 index.
  \item Variables: open, high, low, close, and volume. (OHLCV features)
  \item Sample period: January 1, 2018 to December 31, 2025.
\end{itemize}
\end{frame}

\begin{frame}{Preprocessing}
\begin{itemize}
  \item Address ticker changes, delistings, and newly listed stocks.
  \item Remove invalid observations and holiday-related missing rows.
  \item Drop observations with missing OHLCV information after cleaning.
  \item Final scale: about 2{,}009 trading days and 2{,}520 raw features. \todo{Update this after complete code}
\end{itemize}
\end{frame}

\subsection{Feature Engineering}
\begin{frame}{Target and Features}
\begin{itemize}
  \item Target: binary indicator for next-day S\&P 500 direction.
  \item Added predictors: lags, daily percent changes, ranges, volatility, rolling z-scores, channel position, rolling VWAP, and day-of-week effects.
  \item Raw OHLC levels are removed after feature engineering because they are non-stationary. \todo{Might remove later}
\end{itemize}
\end{frame}

\subsection{Exploratory Data Analysis}

\begin{frame}{Exploratory Data Analysis}
\begin{figure}[h]
    \centering
    \safeincludegraphics[width=0.45\linewidth]{../output/8yrs_sector_correlation.png}
    \safeincludegraphics[width=0.45\linewidth]{../output/8yrs_sector_sp500_correlation_mcap_weighted.png}
    \caption{Correlation across sectors and that of each sector with the S\&P 500 daily return. Highlighted are the top three sectors by absolute market-cap-weighted correlation with the S\&P 500: Healthcare ($|r|=0.1244$), Financial Services ($|r|=0.1198$), and Consumer Defensive ($|r|=0.1112$).}
    \label{fig:sector-correlation}
\end{figure}
\end{frame}

\begin{frame}{Exploratory Data Analysis}
\begin{figure}[h]
    \centering
    \safeincludegraphics[width=0.32\linewidth]{../output/8yrs_sector_Financial_Services_clustered.png}
    \safeincludegraphics[width=0.32\linewidth]{../output/8yrs_sector_Healthcare_clustered.png}
    \safeincludegraphics[width=0.32\linewidth]{../output/8yrs_sector_Consumer_Defensive_clustered.png}
    \caption{Intra-sector correlation matrices for the top three sectors by absolute market-cap-weighted correlation with S\&P 500 returns. From left to right: Financial Services ($|r|=0.1198$), Healthcare ($|r|=0.1244$), and Consumer Defensive ($|r|=0.1112$). Hierarchical clustering groups stocks with similar return dynamics. Warm colours indicate strong positive co-movement while cool colours indicate divergence.}
    \label{fig:top-sectors}
\end{figure}
\end{frame}

\begin{frame}{Exploratory Data Analysis}
  \begin{figure}[h]
    \centering
    \safeincludegraphics[width=0.45\linewidth]{../output/8yrs_sector_volatility_over_time.png}
    \safeincludegraphics[width=0.45\linewidth]{../output/8yrs_smoothed_sector_volatility.png}
    \caption{Sector volatility. Left panel: 21-day-window average volatility. Right panel: its smoothed version by moving average over 40 days.}
    \label{fig:sector-vol}
\end{figure}
\end{frame}

\begin{frame}[t]{Exploratory Data Analysis: Descriptive Statistics}
  \tiny
  \safeinput{../output/8yrs_descriptive_stats.tex}
\end{frame}

\begin{frame}{Exploratory Data Analysis: Descriptive Stats Summary}
  \begin{itemize}
    \item Average daily returns are generally positive, except in 2022.
    \item Return volatility rises sharply in 2020 and again in 2022; Technology has the highest dispersion, while Consumer Defensive has the lowest.
    \item The S\&P 500 up-day proportion is above 0.50 in most years but falls below 0.50 in 2022, indicating mild response imbalance.
  \end{itemize}
\end{frame}

\begin{frame}{Exploratory Data Analysis}
  \textbf{Key findings from EDA:}
  \begin{itemize}
    \item Mild response imbalance.
    \item Certain sectors exhibit stronger correlations with the S\&P 500 returns
    \item Certain segments within these sectors may be driving their overall correlation with the broader market, indicating a high degree of redundancy in data
    \item Different volatility levels suggest that there may be important information contained in the volatility of individual stocks 
    \item A similar time trend in some sector volatility indicates potential data redundancy
  \end{itemize}
\end{frame}




\section{Methods}



\begin{frame}[t,shrink=8]{Method Overview: Data Challenges and solutions}

\begin{block}{High-Dimensional Predictors ($p$ large relative to $n$)}
\begin{itemize}
    \item \textbf{Principal Component Analysis (PCA):} reduce predictors to M principal components (Ch. 14, Section 14.5.1, ESL)
    % ESL reference:
    % Chapter 14, "Unsupervised Learning", Section 14.5.1 "Principal Components", p. 534.
    % The notation below is a slide-friendly score representation:
    % z_{im} = \sum_{j=1}^p x_{ij}\phi_{jm},
    % where \phi_m is the mth loading direction.
    % This is consistent with the principal-component construction in ESL Ch. 14, §14.5.1.
    \[
    z_{im} = \sum_{j=1}^p x_{ij} \phi_{jm}, \qquad m=1,\dots,M
    \]
    where $\phi_m$ is the largest $m$th loading direction.
    
    \item \textbf{Regularization: LASSO} (Ch. 3, Section 3.4.2, ESL)
    % ESL reference for lasso:
    % Chapter 3, Section 3.4.2 "The Lasso", starts p. 68.
    % More generally, ESL Eq. (3.53) in §3.4.3 gives the bridge penalty, p. 72;
    % q=1 corresponds to the lasso and q=2 to ridge.
    % The criterion below is the lasso specialization used on the slide.
    \[
    \hat{\beta}^{\mathrm{lasso}}
    = \arg\min_{\beta}\left\{
    \sum_{i=1}^n (y_i-\beta_0-x_i^\top \beta)^2
    + \lambda \sum_{j=1}^p |\beta_j|
    \right\}
    \]
\end{itemize}
\end{block}
\vfill
{\tiny \textit{Source:} Hastie, Tibshirani, and Friedman (2009), \textit{The Elements of Statistical Learning}, Ch.~3, Ch.~7, and Ch.~14.}

\end{frame}


\begin{frame}[t,shrink=8]{Method Overview: Data Challenges and solutions}
\begin{block}{Multicollinearity}
\begin{itemize}
    \item  Regularization: Ridge and LASSO % ESL reference for ridge:
    % Chapter 3, "Linear Methods for Regression", Section 3.4.1 "Ridge Regression", starts p. 61.
    % Penalized least-squares ridge criterion appears as Eq. (3.43), p. 64.
    \[
    \hat{\beta}^{\mathrm{ridge}}
    = \arg\min_{\beta}\left\{
    \sum_{i=1}^n (y_i-\beta_0-x_i^\top \beta)^2
    + \lambda \sum_{j=1}^p \beta_j^2
    \right\}
    \]
    \item Sector-level feature aggregation
    \item PCA 
    % ESL reference for PCA:
    % Chapter 14, Section 14.5.1 "Principal Components", p. 534.
\end{itemize}
\end{block}

\begin{alertblock}{Time-Series Structure}
\begin{itemize}
    \item Lagged predictors preserve temporal ordering
    \item Models with inherent time-series nature
    %\item Walk-forward validation respects time dependence
    % ESL reference:
    % Chapter 7, "Model Assessment and Selection", starts p. 219.
    % k-fold cross-validation is discussed in this chapter; walk-forward validation
    % is a time-series adaptation of the same model-assessment principle.
\end{itemize}
\end{alertblock}

\vfill
{\tiny \textit{Source:} Hastie, Tibshirani, and Friedman (2009), \textit{The Elements of Statistical Learning}, Ch.~3, Ch.~7, and Ch.~14.}

\end{frame}




\begin{frame}[t,shrink=8]{Method Overview: Models and Model Selection}

\begin{block}{Classification Models}
\begin{itemize}
    \item \textbf{Logistic Regression}
    % ESL reference:
    % Chapter 4, Section 4.4, p. 119.
    % ESL Eq. (4.17) gives the multiclass log-odds formulation. Use binary logistic special case rather than original formula in the book
    \[
    \Pr(G=1\mid X=x)
    =
    \frac{\exp(\beta_0+x^\top\beta)}
    {1+\exp(\beta_0+x^\top\beta)}
    \]
    Easy to be applied with Ridge and LASSO
    \item \textbf{Random Forest}: No explicit restriction for large number of features.
    \begin{itemize}
        \item Ensemble of classification trees
        \item Built on bootstrap samples
        \item Final prediction by majority vote
        % ESL reference:
        % Chapter 15, "Random Forests".
        % Use as a conceptual chapter citation rather than a specific displayed formula.
    \end{itemize}
    
    \item \textbf{Support Vector Machine}: No restriction of number of features, non-linear separable
    % ESL reference:
    % Chapter 12, Section 12.2 p. 417. for SVC and Eqs. (12.7)--(12.8), pp. 420--421. for SVM
    \[
    \min_{\beta,\beta_0,\xi}\ 
    \frac{1}{2}\|\beta\|^2 + C\sum_{i=1}^N \xi_i
    \]
    \[
    \text{subject to } \quad
    y_i(x_i^\top\beta+\beta_0)\ge 1-\xi_i,\qquad
    \xi_i\ge 0,\quad i=1,\dots,N
    \]

    \item Experiment with some neural net models (CNN-2D, LSTM, RNN) 
    
\end{itemize}
\end{block}

\vfill
{\tiny \textit{Source:} Hastie, Tibshirani, and Friedman (2009), \textit{The Elements of Statistical Learning}, Ch.~4, Ch.~7, Ch.~12, and Ch.~15.}

\end{frame}


\begin{frame}[t,shrink=8]{Method Overview: Models and Model Selection}
\begin{block}{Model Tuning and Selection} (Mainly Ch. 7, ESL)
\begin{itemize}
    \item Validation methods: \textbf{Time series} cross-validation
    \item Parameters tuning: selected by 1SE rule
    \item Evaluation metrics:
    \begin{itemize}
      \item Balanced accuracy in parameter tunning (reduced degenerated one)
      \item Between different models: 
      \begin{itemize}
        \item Test accuracy, ROC-AUC, Specificity, Recall (sensitivity), MCC
        \item Diagnostic plots with CV balanced accuracy
      \end{itemize}
    \end{itemize}
\end{itemize}
\end{block}

\vfill
{\tiny \textit{Source:} Hastie, Tibshirani, and Friedman (2009), \textit{The Elements of Statistical Learning}, Ch.~4, Ch.~7, Ch.~12, and Ch.~15.}

\end{frame}

%%%%%%%%%%%%%%%%%%%%%%%%%%%%%%%%%%%%%%%% MAJOR REVISION COMPLETED UP TO HERE%%%%%% LATER INCLUDED NEW RESULTS WITH ANALYSIS BUT IS MINIMALLY REVISED,%%%%%%%%%%%%%%%%%%%%%%%% 


\section{Results}
\subsection{Baseline Models with raw features}
\subsubsection{Logistic Regression Models}
\begin{frame}[t]{Results}
\textbf{Logistic Regression Models}\\
\begingroup
\scriptsize
\setlength{\tabcolsep}{2pt}
\renewcommand{\arraystretch}{0.9}
\safeinput{../output/8yrs_1SE_base_logistic_comparison.tex}
\endgroup
%\FloatBarrier
\end{frame}

\begin{frame}{Results}
\begin{itemize}
  \item Overall accuracy around 50-53\%, slightly better than random guessing
  \item Best test accuracy: LASSO (with or without day of the week) = 0.5572
  \item Best ROC-AUC: Base+PCA: Model without regularization: 0.530
  \item Several combination of Ridge/LASSO with PCA are degenerated
  \item Regularization improve accuracy by about 2-3 pct point, and might reduce predictive variance
  \item Best candidate: 
\end{itemize}
\end{frame}

\begin{frame}{Logistic - Base+PCA}
\begin{figure}[h]
    \centering
    \safeincludegraphics[width=\linewidth]{../output/8yrs_1SE_base_logistic_bias_variance.png}
    \caption{Bias-variance tradeoff for the Base+PCA logistic model, with over/underfitting analysis}
    \label{fig:vb_log_base}
\end{figure}
\end{frame}


\subsubsection{Random Forest Models}
\begin{frame}[t,shrink=25]{Results}
\textbf{Random Forest Models}\\
\safeinput{../output/8yrs_1SE_base_rf_comparison.tex}
%\FloatBarrier
\end{frame}

\begin{frame}{Results}
\begin{itemize}
  \item Raw RF outperforms logistic regression
  \item Applying PCA and LASSO before the random forest worsen test accuracy. 
  \item Simple RF with DOW is much better than without DOW
\end{itemize}
\end{frame}



\begin{frame}{RF with raw features and DOW}
\begin{figure}[h]
    \centering
    \safeincludegraphics[width=\linewidth]{../output/8yrs_base_rf_best_tree.png}
    \caption{RF Tree Classification}
    \label{fig:rf_base_besttree}
\end{figure}
\end{frame}

\subsubsection{SVM Models}
\begin{frame}[t,shrink=25]{Results}
\textbf{SVM Models}\\
\safeinput{../output/8yrs_base_svm_comparison.tex}
%\FloatBarrier
\end{frame}

\subsubsection{NN Models}
\begin{frame}[t,shrink=25]{Results}
\textbf{Neural Net Models}\\
\safeinput{../output/8yrs_base_nn_comparison.tex}
%\FloatBarrier
\end{frame}

\begin{frame}{Results}
\begin{itemize}
  \item Using the RBF kernel improves the test accuracy slightly compared to the linear kernel 
  \item SVM Model using the polynomial kernel is better in ROC-AUC but worse test accuracy overall
  \item LSTM is closer to time-series setting but the raw RNN model is the best performer
\end{itemize}
\end{frame}


\subsection{Extended Feature Models}
\begin{frame}[t,shrink=25]{Results}
\textbf{Logistic Regression}\\
\safeinput{../output/8yrs_logistic_regression_comparison.tex}
\end{frame}

\begin{frame}[t,shrink=25]{Results}
  \textbf{Random Forest}\\
\safeinput{../output/8yrs_random_forest_comparison.tex}
\end{frame}

\begin{frame}[t,shrink=25]{Results}
  \textbf{SVM}\\
\safeinput{../output/8yrs_svm_comparison.tex}
\end{frame}

\begin{frame}[t,shrink=25]{Results}
  \textbf{Neural Network}\\
\safeinput{../output/8yrs_nn_comparison.tex}
\end{frame}

\begin{frame}{Results}
\begin{itemize}
  \item Using engineered features in forcast-
ing does not noticeably improve test accuracy performance
\item Engineered features do not really capture the non-
linear relationship
\item Best performer: RF applied to raw data with DOW
dummies
\item RF variants using dimensionality reduction have lower test accuracy and higher variability.
\end{itemize}
\end{frame}



\begin{frame}{Best performer: RF with raw and DOW}
\begin{figure}[htbp]
  \centering
  \safeincludegraphics[width=0.32\textwidth]{../output/8yrs_base_rf_best_model_bias_variance.png}
   \safeincludegraphics[width=0.32\textwidth]{../output/8yrs_base_rf_best_model_bias_variance_classifier_max_features.png}
    \safeincludegraphics[width=0.32\textwidth]{../output/8yrs_base_rf_best_model_bias_variance_classifier_n_estimators.png}
  \caption{Bias Variance Tradeoff during parameter tunning process for Random Forest with raw data and DOW}
  \label{fig:validation-best-bv}
\end{figure}
\end{frame}


\begin{frame}{Best performer: RF with raw and DOW}
\begin{figure}[htbp]
  \centering
  \safeincludegraphics[width=0.32\textwidth]{../output/8yrs_base_rf_best_model_train_test.png}
   \safeincludegraphics[width=0.32\textwidth]{../output/8yrs_base_rf_best_model_train_test_classifier_max_features.png}
    \safeincludegraphics[width=0.32\textwidth]{../output/8yrs_base_rf_best_model_train_test_classifier_n_estimators.png}
  \caption{Over/underfitting during parameter tunning process for Random Forest with raw data and DOW}
  \label{fig:validation-best-bv}
\end{figure}
\end{frame}


\begin{frame}{Results}
  \textbf{Complexity}
\begin{figure}[htbp]
  \centering
  \safeincludegraphics[width=0.95\textwidth]{../output/8yrs_family_leaders_accuracy_vs_cv_sd.png}
  \caption{Accuracy and Proxy for Noise of The Best Models within Model Families and Data Set Performer}
  \label{fig:complexity2}
\end{figure}
\end{frame}


\begin{frame}{Results}
  \textbf{Complexity}
\begin{figure}[htbp]
  \centering
\safeincludegraphics[width=0.95\textwidth]{../output/8yrs_family_leaders_cv_accuracy_vs_cv_sd.png}
  \caption{Cross-Validation: Balanced Accuracy for The Best Candidate within Model Families and Data Set}
  \label{fig:complexity1}
\end{figure}
\end{frame}




\begin{frame}{Results}
  \textbf{Prediction Visualization}
\begin{figure}[htbp]
  \centering
  \safeincludegraphics[width=0.95\textwidth]{../output/8yrs_prediction_distribution_base_random_forest_py_rf_dow.png}
  \caption{Predicted Probability Distribution by True Class on The Test Set}
  \label{fig:pred-dist-base-rf-dow}
\end{figure}
\end{frame}

\begin{frame}{Results}
\begin{itemize}
  \item RF using raw data with DOW is still a good candidate in diagnostic plots
  \item LSTM is the second best
  \item LASSO regularization (L1) demonstrates superior highest accuracy (0.5486) and the only ROC-AUC above 0.50 (0.5342)
  \item  SVM with RBF kernel exhibits degenerate behavior, predicting all samples as ``Up'' 
  \item The dimensionality reduction models (PCA, LASSO-reduced, Ridge-reduced) show moderate performance, potentially due to losing important predictive signals. 
\end{itemize}
\end{frame}


\begin{frame}{Results - Summary}
\begin{itemize}
  \item Predictive performance of RF is the best
  \item More advanced models suffer from higher variablity (reflected in Validation SE) while not always give better performance
  \item RF models suggest that the EMH is not supported.
  \item Prediction visualisation suggest further tuning of decision threshold.
\end{itemize}
\end{frame}

\section{Limitations \& next steps}
\begin{frame}{Limitation}

  \begin{itemize}
    \item Data: 8 years rather than longer, causing mild response bias
    \item Models: Limited grid-search lists, too simple neural net architectures.
    \item Evaluation: Hard threshold decision of 0.5
    \item Time constraints: limited time for extensive hyperparameter tuning and model decision threshold
    \item Computational resources: limited in training more complex models
    \end{itemize}
\end{frame}

\begin{frame}{Next steps}

\begin{itemize}
  \item More extensive list for hyperparameter tuning and various choices of model decision threshold
  \item Possibly experiment with cSE rule rather than 1SE rule.
  \item Varying threshold of decision for classifier
  \item Include longer time series, shorter time window
  \item Explore more sophisticated models (more complicated deep learning architectures)
  \item Potential of combination with economic method of time series forecasting (e.g., ARIMA, GARCH)
\end{itemize}
\end{frame}

\section{Final summary}
\begin{frame}{Summary}

\begin{itemize}
  \item Supporting evidence for challenging EMH
  \item Relying solely on historical daily price and volume data may not be sufficient for good performance
  \item Even not very good, result still can be employed in trading strategy.
  \item Future work : more extensive parameter grids, varying decision threshold, more models exploration, more comprehensive evaluation framework
  \end{itemize}
\end{frame}


% --- References ---
\section{References}
\begin{frame}[allowframebreaks]{References}
\nocite{*}
\bibliographystyle{apacite}
\bibliography{Bibliography}
\end{frame}
\end{document}
